% =========================================================
%  Section : Le théorème de Donaldson
% =========================================================

\subsection{Donaldson(1983)}

Le théorème de Donaldson, publié en 1983, est une obstruction à la lissabilité d'une
$4$-variété topologique qui s'applique sans aucune hypothèse de structure spin. Il concerne
les formes d'intersection définies positives et impose une contrainte de nature
arithmétique~: toute forme définie positive réalisable par une variété lisse doit être
diagonalisable sur $\mathbb{Z}$.

\begin{theorem}[Donaldson, 1983]\label{thm:donaldson}
  Soit $M$ une $4$-variété lisse compacte orientée fermée simplement connexe. Si la forme
  d'intersection $Q_M$ est définie positive, alors $Q_M$ est isomorphe à
  $\langle 1 \rangle^{\oplus n}$ sur $\mathbb{Z}$, où $n = b_2(M)$.
\end{theorem}

Autrement dit, dans la base appropriée, la matrice de $Q_M$ est la matrice identité. Toute
forme définie positive qui n'est pas diagonalisable sur $\mathbb{Z}$ est donc inaccessible
dans la catégorie lisse, quelle que soit la structure spin de la variété.

\medskip

Ce résultat contraste fortement avec la situation dans la catégorie topologique~: le théorème
de Freedman, présenté dans la section suivante, garantit que toute forme unimodulaire est
réalisable par une $4$-variété topologique. Le théorème de Donaldson révèle ainsi un écart
brutal entre les deux catégories, qui est un phénomène propre à la dimension $4$.

\subsubsection*{Conséquences}

Le théorème de Donaldson classe entièrement les formes définies positives réalisables par des
variétés lisses~:

\begin{corollary}\label{cor:donaldson_classification}
  Une forme bilinéaire symétrique entière définie positive est la forme d'intersection d'une
  $4$-variété lisse compacte orientée fermée simplement connexe si et seulement si elle est
  isomorphe à $\langle 1 \rangle^{\oplus n}$ pour un certain $n \geq 0$.
\end{corollary}

Combiné au théorème de Freedman, ce résultat produit une famille d'exemples de variétés
topologiques sans structure lisse~: toute forme définie positive unimodulaire non
diagonalisable sur $\mathbb{Z}$ est réalisable topologiquement mais pas lisément. L'exemple
le plus économique est la forme $E_8$, à laquelle la section suivante est consacrée.

\medskip

On note enfin que le théorème de Donaldson englobe le théorème de Rokhlin dans le cas défini
positif~: si $Q_M$ est définie positive et lisse, alors $Q_M \cong \langle 1 \rangle^{\oplus
n}$, qui est une forme impaire. En particulier $M$ n'est pas spin, et Rokhlin ne
s'applique pas. Les deux théorèmes sont donc complémentaires plutôt que redondants~: Rokhlin
contraint les formes paires via la divisibilité de la signature, Donaldson exclut toutes les
formes définies positives non diagonales indépendamment de la parité.